\documentclass[twocolumn]{autart}  





\usepackage{graphicx}          
\usepackage{amsmath}
\usepackage{amssymb}
\usepackage{amsfonts}
%\usepackage{amsthm}
\usepackage[dvipsnames]{xcolor}
%\usepackage{hyperref}
\usepackage[hypertexnames=false]{hyperref}


\usepackage{comment}

\newcommand{\B}[1]{{\textcolor{blue}{[bruno: \textit{#1}]}}}
\newcommand{\LS}[1]{{\textcolor{ForestGreen}{[LS: \textit{#1}]}}}
\newcommand{\LP}[1]{{\textcolor{purple}{[LP: \textit{#1}]}}}
%\usepackage{boldsymbol}

\newcommand{\Luc}{\color{magenta} }
\newcommand{\ludo}{\color{ForestGreen}}






%%%%%%%%%%%%%%%%%%%%%%%%%%%%%%%%%%%%%%%%%%%%%%%%%%%%%%%%%%%%%%%%%%%%
%%%%%%%%%%%%%%%%%%%%%%%%%%%%%%%%%%%%%%%%%%%%%%%%%%%%%%%%%%%%%%%%%%%%
%%%%%%%%%%%%%%%%%%%%%%%%%%%%%%%%%%%%%%%%%%%%%%%%%%%%%%%%%%%%%%%%%%%%


\newtheorem{theorem}{Theorem}[section]
\newtheorem{corollary}[theorem]{Corollary}
\newtheorem{proposition}[theorem]{Proposition}
\newtheorem{lemma}[theorem]{Lemma}

\theoremstyle{definition}
\newtheorem{definition}[theorem]{Definition}
\newtheorem{assumption}[theorem]{Assumption}

\theoremstyle{remark}
\newtheorem{remark}[theorem]{Remark}
\newtheorem{example}[theorem]{Example}


%%%%%%%%%%%%%%%%%%%%%%%%%%%%%%%%%%%%%%%%%%%%%%%%%%%%%%%%%%%%%%%%%%%%
%%%%%%%%%%%%%%%%%%%%%%%%%%%%%%%%%%%%%%%%%%%%%%%%%%%%%%%%%%%%%%%%%%%%
%%%%%%%%%%%%%%%%%%%%%%%%%%%%%%%%%%%%%%%%%%%%%%%%%%%%%%%%%%%%%%%%%%%%

\newcommand{\R}{\mathbb{R}}
\newcommand{\N}{\mathbb{N}}
\newcommand{\diff}{\mathrm{d}}
\newcommand{\id}{\mathrm{id}}
\newcommand{\Cov}{\mathrm{Cov}}
\newcommand{\Conv}{\mathrm{Conv}}
\newcommand{\supp}{\mathrm{supp}}
\newcommand{\Tr}{\mathrm{Tr}}


\newcommand{\U}{\mathcal{U}}
\newcommand{\V}{\mathcal{V}}
\renewcommand{\P}{\mathcal{P}}

\newcommand{\yb}{\mathbf{y}}
\newcommand{\ub}{\mathbf{u}}
\newcommand{\Jb}{\mathbf{J}}
\newcommand{\Mb}{\mathbf{M}}
\newcommand{\Sb}{\mathbf{S}}
\newcommand{\diag}{\mathrm{diag}}
\renewcommand{\S}{\mathcal{S}}

\newcommand{\etab}{\boldsymbol{\eta}}
\newcommand{\Sigmab}{\boldsymbol{\Sigma}}
\begin{document}

\begin{frontmatter}


\title{Optimal experimental design for parameter estimation in linear controlled stochastic system with application to retinal neural networks}

\thanks[footnoteinfo]{This paper was partially supported by ANR...}

\author[inria]{Ludovic Sacchelli}\ead{ludovic.sacchelli@inria.fr},    % Add the 
\author[inria2]{Bruno Cessac}\ead{bruno.cessac@inira.fr},               % e-mail address 
\author[I3S]{Luc Pronzato}\ead{luc.pronzato@i3s.unice.fr}  

\address[inria]{Inria , Universit\'e C\^ote d'Azur, CNRS, LJAD, France}  
\address[inria2]{Inria}  
\address[I3S]{I3S}            


          
\begin{keyword}                           
Optimal experimental design                                      
\end{keyword}                             
                                          
                                          


\begin{abstract}                         
We discuss the problem of experiment design for continupous time linear controlled stochastic systems. Such systems can be used to model the neural network of the retina and help design experiments meant to estimate parameters in the retinal response.
\end{abstract}

\end{frontmatter}

\section{Introduction}


Contemporary technology in neuroscience allow for experiments to be conducted in “closed-loop”. One application concerns the characterisation of the retina. This membrane, which lines the back of the eye and is approximately half a millimeter thick (in humans), is capable of converting photons from the outside world into trains of impulses (action potentials) that are transmitted to the brain via the optic nerve. To better understand how the structure and dynamics of this organ enable it to optimally encode visual scenes, it is subjected to different types of stimuli and the response is recorded, so that mathematical models can then be developed. In closed loop experiments, the retina is subjected to simple stimuli (such as “white noise”), and the response is recorded, which is then used to adjust the parameters of a preliminary model. Based on this model, a new stimulus is generated using an optimality criterion. This new stimulus is then subjected to the retina, and so on. \B{Ajouter refs.}

However, this approach encounters several problems. Firstly, the experiments must be conducted under stringent time constraints. This means that measurements, parameter estimates/model adjustments, and stimulus determinations must be performed over short periods of time, with the aim of obtaining at least ten loops per experiment. Second, the models have a large dimensionality. In fact, closed-loop experiments typically aim to determine the properties of individual neurons in the retina. But more recently, as it is the case here, researchers wanted to determine the collective properties of a large number of cells (several hundred) interacting in a network. This implies high dimensionality in phase space. The parameter space can initially be reduced by using simplified phenomenological models \cite{kartsaki-hilgen-etal:24}, as we use here. Finally, optimality criteria are usually defined via the characteristics of individual cells, for example, maximum activity achieved for specific "preferred" stimuli. We would like to develop optimality criteria that take into account the \textit{collective dynamics} of retinal cells interacting via a synaptic network in response to non-stationary stimuli. Furthermore, this criterion should incorporate a notion of optimality in the way visual \textit{information} is encoded.

In this article, we address these questions from the natural perspective of optimal experiment design (OED) theory, within the linear version of the framework of a retinal model proposed in \cite{kartsaki-hilgen-etal:24}. In this model, neural activity in the retina is described by a linear continuous time system subjected to non-stationary and non-homogeneous inputs and driven by white noise, that is then measured at its base. For this reason the question is to apply the framework of OED to linear (multi-input-multi-output) controlled stochastic systems (e.g. multi-dimensional Ornstein--Uhlenbeck processes).

\LS{Write here something about the theory of OED. Write also an up-to-date discussion of what exists for OU processes.}

In the present paper, we formulate the problem of parameter estimation and stimulus optimality in this context, adopting the maximum-likelihood approach of [...]. By restricting the set of controls to piecewise constant (in time) inputs that update their values in synchronization with the measurement times, it is possible to create highly structured and correlated time-series to be analyzed.
The Fisher information matrix plays a central role in this method, and we show that, in this model, it can be calculated explicitly, along with the first and second derivatives of the maximum likelihood energy.
The Fisher information matrix naturally arises in the parameter estimation problem, via a central limit theorem, but also in the determination of optimal stimuli. Indeed, potential optimality criteria are defined as functions of this matrix \cite{Pukelsheim2006,Fedorov72}. We study the mathematical properties of a criterion based on the determinant of this matrix, properties that are particularly relevant to the numerical problem of stimulus design. We also discuss how such optimality criteria, explicitly involving spatio-temporal correlations induced by collective dynamics (called “noise correlations” in the neuroscience literature), may offer better options than those currently used in the study of the retina by closed-loop stimulation.

In the paper, we first discuss the establishment of the theory of OED in application and Fisher matrix optimization in the context of linear controlled stochastic systems, then in a second time we will focus on the description of the proper experimental model and numerical simulations of the method.

\section{Setting}


 
\subsection{Dynamics}
\label{S:dynamics}

Let us start by introducing the linear dynamical model. Let $n,m,p,q\in \N$. The vector of parameter of the system is denoted by $\theta\in \Theta\subset \R^p$, with $\Theta$ the set of admissible parameters assumed to be compact. The dependence of the system to the parameter is assumed to enter through the mapping $A:\Theta\to \R^{n\times n}$, that we assume to be a mapping of regularity $C^1$ at least. In the applications, we essentially assume $A$ to be affine. Let $B\in \R^{n\times m}$, $C\in \R^{q\times n}$, and $D\in \R^{n\times n}$. Let $W$ be a standard $n$-dimensional Brownian motion. For a given $T>0$, and $u\in L^\infty([0,T],\R^m)$, for a given parameter $\theta\in \Theta$,  we consider the controlled stochastic differential equation:
\begin{equation}\label{E:sde1}
\diff x(t)=A(\theta) x(t) + B u(t)+D \diff W(t).
\end{equation}



\begin{assumption}\label{As:Gap}
We assume there exists $\gamma>0$ such that for all $\theta\in \Theta$, for any eigenvalue $\lambda$ of $A(\theta)$, we have
%For all $\theta\in \Theta$, $A(\theta)$ is Hurwitz, that is, all eigenvalues of $A(\theta)$ have negative real part. Furthermore, we assume there exists $\gamma>0$ such that for any eigenvalue $\lambda$ of $A(\theta)$ with any $\theta\in \Theta$,
$$
\mathrm{Re}(\lambda)<-\gamma.
$$
\end{assumption}

For a given deterministic initial condition $x(0)=x_0\in \R^n$, \eqref{E:sde1} admits a unique strong solution  (see, e.g., \cite[Theorem 5.2.1]{oksendal2013stochastic}).
{\ludo
Then assume that we let evolve \eqref{E:sde1} under the null input with some initial condition $x_0\in \R^n$. As $t\to \infty$, the distribution of $x(t)$ converges towards an invariant distribution $x_\infty\sim \mathcal{N}(0,\hat\Sigma)$, where the covariance matrix $\hat \Sigma$ is the unique symmetric positive definite solution of the Lyapunov equation $A\hat \Sigma+\hat \Sigma A^\top + DD^\top = 0$.

As a consequence, we make the assumption that each experiment starts with the system at rest so that $x_0$ is not deterministic but rather assumed to follow the invariant distribution $x_0\sim \mathcal{N}(0,\hat\Sigma)$. Under these conditions, \eqref{E:sde1} still admits a unique strong solution (following from the same reference),}
given by, for any $t\in [0,T]$,
\begin{multline*}%\label{E:VarConst}
x(t)=e^{A(\theta)t}x_0+\int_0^t e^{A(\theta)(t-s)}Bu(s)\diff s
\\+
\int_0^t e^{A(\theta)(t-s)}D \diff W(s).
\end{multline*}
At any time $t\in [0,T]$, the \emph{output} of the system is given by $y(t)=C x(t)$.


Let $N\in \N$, we set $\Delta=T/N$ and $t_k=k \Delta$, for $k=0,\dots, N$. We also denote $x_k=x(t_k)$. If we compute increments of $x$ between two dates $t_{k},t_{k+1}$, $0\leq k<N$, we get:
\begin{multline}\label{E:rec1}
x_{k+1}=e^{\Delta A(\theta)}x_k+\int_{t_k}^{t_{k+1}}e^{A(\theta)(t_{k+1}-s)}Bu(s)\diff s
\\
+
\int_{t_k}^{t_{k+1}}e^{A(\theta)(t_{k+1}-s)}D \diff W(s).
\end{multline}
\begin{assumption}\label{A:periodic}
Any admissible input $u$ is such that there exists a % $\Delta$-periodic, essentially-bounded function $U:\R\to \R^{m\times m}$, and 
a sequence $(u_k)_{0\leq k<N}\in \left(\R^m\right)^N$ such that: 
\begin{equation*}
    u(t)=u_k, \quad  t\in [t_k,t_{k+1}), 0\leq k< N.
\end{equation*}
We denote by $\U$ the set of these inputs. 
\end{assumption}
Set $F=e^{\Delta A(\theta)}$,
$
G=\int_{0}^{\Delta}e^{A(\theta)(\Delta-s)}BU(s)\diff s,$
$W_k=\int_{t_k}^{t_{k+1}}e^{A(\theta)(t_{k+1}-s)}D \diff W(s),$
%
where for simplicity we omit the dependence in $\theta$.
Equation~\eqref{E:rec1} becomes 
\begin{equation*}
x_{k+1}=F x_k + G u_k + W_k,
\end{equation*}
with $(W_k)_{0\leq k<N}$ an i.i.d. sequence of centered Gaussian random vectors with covariance matrix $\Sigma$, i.e.\ $W_k\sim \mathcal{N}(0,\Sigma)$, with
$$
\Sigma = \int_0^\Delta e^{A(\theta)s}D D^\top  e^{A^\top(\theta)s}\diff s.
$$


% \begin{remark}\LP{C'est déjà dit quelques lignes plus haut\ldots}
% The matrices $F\in \R^{n\times n}$, $G\in \R^{n\times m}$ and $\Sigma\in \R^{n\times n}$ all implicitly depend on $\theta$. Unless necessary, we do not mention this fact in the following.
% \end{remark}
\begin{remark}
The integral definitions of $G$ and $\Sigma$ can be  impractical from a computational standpoint as they use an exponentiation of $A(\theta)$. Hence we may prefer definitions based on ordinary differential equations: $G=\tilde G(\Delta)$ and $\Sigma=\tilde \Sigma(\Delta)$, where 
$$
\dot{\tilde{G}}(t)= A(\theta)\tilde{G}(t)+BU(t),\quad {\tilde{G}}(0)=0\in \R^{n\times m},
$$
$$
\dot{\tilde{\Sigma}}(t)= A(\theta)\tilde{\Sigma}(t)+\tilde{\Sigma}(t)A(\theta)^{\top}+D D^\top ,\,{\tilde{\Sigma}}(0)=0\in \R^{n\times n}.
$$
\end{remark}
Finally, we assume there exists $(\varepsilon_k)_{k< 0 \leq N}$, i.i.d.\ centered Gaussian random vectors with covariance matrix $\Sigma'\in \R^{q\times q}$, that is $\varepsilon_k\sim \mathcal{N}(0,\Sigma')$; we assume $\Sigma'$ is  positive-definite. A measurement  at $t_k$, $0< k \leq  N$, thus follows
\begin{equation}\label{E:measurement}
y_k(u)
=\eta(t_k,\theta,u)+CF^{k}x_0+\sum_{l=0}^{k-1} C F^{k-l-1}W_{l}+\varepsilon_k,
\end{equation}
with $\eta$ denoting the deterministic output
\begin{equation}\label{E:output}
\eta(t_k,\theta,u)
=
\sum_{l=0}^{k-1} C F^{k-l-1}Gu_{l} .
\end{equation}

\begin{remark}
\LS{To be removed:}
In keeping with the uniform dissipativity assumption~\ref{As:Gap}, we will assume $x_0=0$. In all generality $x_0$ may rather be assumed to be probabilistically distributed.
\end{remark}

\subsection{Distribution of the output}



First, let us compute $\Cov(y_k,y_l)$ in all generality. Assume $k\geq l$, since $\Cov(W_i,W_j)=\delta_i^j \Sigma$ ($\delta_i^j$ Kronecker symbol), we get:
$$
\begin{aligned}
\Cov(y_k,y_l)
&=
\sum_{i=0}^{k-1}\sum_{j=0}^{l-1}CF^{k-i-1}\Cov(W_i,W_j)(F^\top)^{l-j-1}C^\top
+\delta_k^l\Sigma'\\
&=
\sum_{i=0}^{l-1}CF^{k-i-1}\Sigma(F^{l-i-1})^\top C^\top
+\delta_k^l\Sigma'
\\
&=
CF^{k-l}\bigg(\sum_{j=0}^{l-1}F^{j}\Sigma(F^j)^{\top}\bigg)C^\top
+\delta_k^l\Sigma'. %\qquad\qquad (j=l-1-i)
\end{aligned}
$$
Letting $\yb(u)\in \R^{Nq}$ denote the vector concatenation of $y_1(u),\dots ,y_N(u)$, $\yb^\top=\begin{pmatrix}
y_1^\top & \cdots & y_N^\top
\end{pmatrix}$,
the full covariance matrix of $y$ is then given by a $Nq\times Nq$ matrix denoted $\Sigmab$. To understand the structure of the blocks in this matrix, it may be easier to consider $\Sigma_l=\sum_{j=0}^{l-1}F^{j}\Sigma(F^j)^{\top}$ (in particular $\Sigma_{l+1}=\Sigma+F \Sigma_l F^\top$ for any $l\geq 1$, with $\Sigma_1=\Sigma$), and to decompose $\Sigmab$ into
$$
\Sb
=
\left(\Cov(x_k,  x_l)\right)_{1\leq k,l\leq N}
=
\begin{pmatrix}
\Sigma_1 & \Sigma_1 F^\top & \Sigma_1 (F^2)^\top & \cdots 
\\
F \Sigma_1 & \Sigma_2  & \Sigma_2 F^\top & \cdots 
\\
F^2 \Sigma_1 & F\Sigma_2  & \Sigma_3  & \cdots 
\\
\vdots&\vdots&\vdots&\ddots
\end{pmatrix},
$$
$$
\Sigmab
=
\diag(C,\dots, C)
\cdot\Sb\cdot
\diag(C,\dots, C)^\top
+
\diag(\Sigma', \dots, \Sigma'),
$$
Here, $\Sigmab$ is implicitly a function of the true parameter $\bar\theta$. %, as the self-covariance matrix of the measured output.
More generally, we write $\Sigmab(\theta)$ for the same construction carried out with an arbitrary parameter $\theta$ (using $F(\theta)$ and $\Sigma(\theta)$).

\begin{lemma}
\ludo TBD: decrease of the amplitude of blocks along the diagonal
\end{lemma}



\subsection{Identifiability}
Identifiability addresses the question of whether $\theta$ can be estimated uniquely under ideal conditions involving a potentially infinite set of experiments. Such ideal conditions would allow us to eliminate the effect of noise on observations, and here the questions boils down to the validity of the following implication: $\eta(t,\theta',u)=\eta(t,\theta,u)$ for all possible $u$ and $t$ implies $\theta'=\theta$. To avoid pathological situations (see the example below), {\em structural global identifiability} requires the implication to be true for almost all $\theta$ in $\Theta$ (in the Lebesgue measure sense on $\Theta$). When the property is only true locally; that is, for almost all $\theta\in\Theta$, there exists a neighborhood $V(\theta)$ such that $\eta(t,\theta',u)=\eta(t,\theta,u)$ for all possible $u$ and $t$, and $\theta'\in V(\theta)$, implies $\theta'=\theta$, the model is {\em structurally locally identifiable}. This condition is sufficient to make the Fisher information matrix (see Section~\ref{S:FIM}) strictly positive definite for almost all $\theta$. One may refer to \cite{KoopmansR50} for an early reference on structural identifiability. For linear time-invariant dynamical systems, a classical approach for testing a model for identifiability relies on the Laplace transform; see, e.g., \cite{BellmanA70}. \cite{WPl93}. 
We illustrate it through an example.


\begin{example}
Take $n=p=2$, $q=m=1$, 
$$
A(\theta)=\begin{pmatrix}\theta_1 -1 & 1-\theta_2 \\
\theta_2 & \theta_1 -1 
\end{pmatrix} \mbox{ and } B=\begin{pmatrix} 0\\
1
\end{pmatrix} \,.
$$
The input-output behavior is entirely characterised by the Laplace transform of the input-output transfer matrix; that is,
$H(s,\theta)=C (s\times \id_{\R^{n\times n}}-A(\theta))^{-1} B$, with $s$ the Laplace variable.
%$H(s,\theta)=C (s I_n-A(\theta))^{-1} B$, with $s$ the Laplace variable and $I_n$ denoting the $n$-dimensional identity matrix.

As here $m=q=1$, when $C=(1\ 0)$ the Laplace transform is
$$
H(s,\theta)= \frac{1-\theta_2}{s^2+2s(1-\theta_1)+(1-\theta_1)^2+\theta_2(\theta_2-1)} \,.
$$
The equality $H(s,\theta')=H(s,\theta)$ for all $s$ implies $\theta'=\theta$ and the model is  structurally globally identifiable (note that for the particular value $\theta_2=1$ the output is identically zero whatever the input $u$, hence the requirement of a structural property being valid for {\em almost all} $\theta$). 
Take now $C=(0\ 1)$. It gives
$$
H(s,\theta)= \frac{s+1-\theta_1}{s^2+2s(1-\theta_1)+(1-\theta_1)^2+\theta_2(\theta_2-1)} \,.
$$
The parameter $\theta_1$ is then structurally globally identifiable, but $H(s,\theta')=H(s,\theta)$ for $\theta'=(\theta_1,1-\theta_2)$ so that $\theta_2$ is only structurally locally identifiable.

\end{example}

In the following we always assume that the model is structurally locally identifiable. Note that this can be quite demanding in terms of conditions set on $B$ and $C$. For instance, if $A=\diag\{\theta_i,\, i=1,\ldots,n\}$, structural global (or even local) identifiability requires that both $B$ and $C$ have  rank $n$. A more precise condition will be given in Section~\ref{S:FIM}, see \eqref{estimability}.



\section{Design of optimal experiments}

%\subsection{Generalized least-squares}
\subsection{Maximum-likelihood estimation}

From now on, we assume that there is no modelling error; that is, that the measured outputs $y(u)$ in \eqref{E:measurement} are generated by a model with a parameter vector denoted $\bar\theta\in \Theta$, considered as the unknown true value of $\theta$ that we wish to estimate.





%(Remark that the $(k,l)$-block of $\Sb$ can also be written $F^{\max(0,k-l)} \Sigma_{\min(k,l)} {(F^\top)}^{\max(0,l-k)}$.)



Denoting $\etab(\theta,u)\in \R^{Nq}$ the vector concatenation of $\eta(t_1,\theta,u),\dots ,\eta(t_N,\theta,u)$, that is  $\etab^\top(\theta,u)=\begin{pmatrix}
\eta(t_1,\theta,u)^\top & \cdots & \eta(t_N,\theta,u)^\top
\end{pmatrix}$,
we can express the 
%
log-likelihood criterion (up to a change of sign) as
\begin{equation}\label{E:GLS1}
J(\theta,u)
=
(\yb(u)-\etab(\theta,u))^\top 
\Sigmab^{-1}(\theta)
(\yb(u)-\etab(\theta,u))+\log\det \Sigmab.
\end{equation}


In the following we assume that the experiment used for  estimating of $\theta$ consists of the successive application of $K$ different inputs $u^{(1)},\dots,u^{(K)} \in \U$, producing $K\times N\times q$ scalar observations in total.
The associated full cost-functional (to be minimised) for the maximum-likelihood estimation of $\theta$ is then 
\begin{equation}\label{E:GLS1.5}
\Jb\left(\theta,\left(u^{(1)},\dots,u^{(K)} \right)\right)
=
\frac{1}{K}\sum_{k=1}^K 
(\yb(u^{(k)})-\etab(\theta,u^{(k)}))^\top 
\Sigmab^{-1}(\theta)
(\yb(u^{(k)})-\etab(\theta,u^{(k)}))+\log\det \Sigmab.
\end{equation}
Note that $\Sigmab$ does not depend on $u$. 

\subsection{Asymptotic estimation via experiment sequences}
%\subsubsection{Fisher information matrix}
\label{S:FIM}
By Assumption~\ref{A:periodic}, $\mathcal{U}$ is in bijection with $\R^N$ \LP{plutôt $\R^{mN}$?} and we can assume $\mathcal{U}$ {to be} endowed with the same topology and Borel $\sigma$-algebra  by push-forward. We denote by $\P(\U)$ the set of probability measures with respect to this $\sigma$-algebra. If $\mu_K\in \P(\U)$ is the discrete uniform probability measure supported on $(u^{(1)},\dots, u^{(K)})$, then we can rewrite $\Jb$ as a function of $\mu_K$:
\begin{equation}\label{E:GLS3}
\Jb(\theta,\mu_K)
=
\int_\U
(\yb(u)-\etab(\theta,u))^\top 
\Sigmab^{-1}(\theta)
(\yb(u)-\etab(\theta,u))\diff \mu_K(u)+\log\det \Sigmab(\theta).
\end{equation}
This definition of $\Jb$ extends  to any probability $\mu\in \P(\U)$. We further assume that we work within a compact subset $\V$ of $\U$ and that all measures $\mu$ under consideration have support in $\V$. We denote by $\P(\V)$ the subset of $\P(\U)$ of probability measures with support $\V$.
Under these conditions, \eqref{E:GLS1.5} can be reinterpreted as an empirical realisation of some $\mu\in \P(\V)$, where the $u^{(k)}$ are independently sampled from $\mu$. We may assume that the input sequences $u$ in $\V$ are somewhat normalised: for instance, 
denoting $\ub^\top=\begin{pmatrix}
u_0^\top & \cdots & u_{N-1}^\top
\end{pmatrix}\in \R^{Nm}$,
we may impose the global power constraint $\frac{1}{Nm} \ub^\top\ub=1$, or power constraints on each component of the $u_k$, i.e., 
$\frac1N \sum_{k=0}^{N-1} \{u_k^2\}_i=1$ for $i=1,\ldots,m$. 


\paragraph{Strong consistency.}
We assume that $\Theta$ is compact and coincides with the closure of its interior, and that $\bar\theta\in\Theta$. We also assume that $A(\theta)$ is continuous on $\Theta$ and that $\eta(t,\theta,u)$ is bounded for all $u\in\V$. Under these conditions, $\etab(\theta,u)$ is bounded on $\Theta\times\V$ and continuous on $\Theta$ for any $u\in\V$; moreover, $\Sigmab(\theta)$ is continuous on $\Theta$, $\Sigmab(\bar\theta)$ is strictly positive definite and $\Sigmab^{-1}(\theta)$ is bounded on $\Theta$. 




For a given probability measure $\mu\in \P(\V)$, let $(u^{(k)})_{k\in \N}$ be a sequence of independent random variables in $\V$ distributed according to $\mu$. We then define $\mu_K=\frac{1}{K}\sum_{k=1}^K \delta_{u^{(k)}}$ the associated empirical realisation of $\mu$ based on the $K$ first samples.
Following the same arguments as in \cite[Th.~3.22]{PP2013}, we obtain that, as $K$ tends to infinity, $\Jb(\theta,\mu_K)$ tends almost surely and uniformly in $\theta$ to 
$$
\bar\Jb(\theta,\mu) = \int_\V (\etab(\theta,u)-\etab(\bar\theta,u))^\top 
\Sigmab^{-1}(\theta)
(\etab(\theta,u)-\etab(\bar\theta,u))\diff \mu(u)+2\,D_{KL}(\theta\|\bar\theta) + k + \log\det \Sigmab(\bar\theta)\,,
$$
where 
$$
D_{KL}(\bar\theta\|\theta)=\frac12\left\{ \Tr\left[\Sigmab^{-1}(\theta)\Sigmab(\bar\theta)\right]-k+\log \frac{\det \Sigmab(\theta)}{\det \Sigmab(\bar\theta)} \right\} \quad (\geq 0)
$$
is the Kullback-Leibler divergence between two normal distributions having the same mean and covariances $\Sigmab(\bar\theta)$ and $\Sigmab(\theta)$. This implies that, under the estimability condition over $\Theta$:
\begin{equation}\label{estimability}
\left[
\int_\V (\etab(\theta,u)-\etab(\bar\theta,u))^\top 
\Sigmab^{-1}(\theta)
(\etab(\theta,u)-\etab(\bar\theta,u))\diff \mu(u) =0 \; \text{ and }\;
D_{KL}(\bar\theta\|\theta)=0 
 \right] \Longrightarrow \theta=\bar\theta,    
\end{equation}
the maximum likelihood estimator $\widehat\theta^K$ that minimises $\Jb(\theta,\mu_K)$
with respect to $\theta\in\Theta$ converges almost surely to $\bar\theta$ as $K\to\infty$.


\paragraph{Asymptotic normality.}
Assuming, moreover, that $\bar\theta\in \mathrm{int}(\Theta)$ and that $A(\theta)$ is twice continuously differentiable with respect to $\theta\in\mathrm{int}(\Theta)$, we obtain that $\widehat\theta^K$ is asymptotically normally distributed (see, e.g., \cite[Th.~3.24]{PP2013}):



%Under classical regularity assumptions\footnote{Check which ones.}  on the functions $\etab(t,\cdot,u)$ and on the admissible set $\Theta$ for $\bar\theta$ ($\Theta$ is compact and $\bar\theta\in \mathrm{int}(\Theta)$), the maximum likelihood estimator $\widehat\theta^K$ that minimises $\Jb(\theta,\mu_K)$ with respect to $\theta\in\Theta$ converges \B{Lebesgue ?} a.s. to $\bar\theta$ when $K\to\infty$ and is asymptotically normal:
\begin{equation}\label{E:TCL}
\sqrt{K} (\widehat\theta^K-\bar\theta) \xrightarrow[K\to +\infty]{d} z \sim \mathcal{N}\left(0,\Mb^{-1}(\bar\theta,\mu)\right),
\end{equation}
where $\Mb(\theta,\mu)\in \R^{p\times p}$ denotes the Fisher information matrix:
\begin{equation}\label{E:defM}
\Mb(\theta,\mu) = 
\int_\V \frac{\partial\etab^\top(\theta,u)}{\partial\theta}
\Sigmab^{-1}(\theta) \frac{\partial\etab(\theta,u)}{\partial\theta^\top}  \diff \mu(u)+Q(\theta),
\end{equation}
with $Q(\theta)$ the symmetric matrix:
\begin{equation}\label{E:defQ}
Q_{ij}(\theta) = \frac{1}{2} \Tr\left[\Sigmab^{-1}(\theta)\frac{\partial\Sigmab(\theta)}{\partial\theta_i}\Sigmab^{-1}(\theta)\frac{\partial\Sigmab(\theta)}{\partial\theta_j}\right], \quad 1\leq i,j\leq p.
\end{equation}
Under these circumstances, optimising $\Mb(\theta,\mu)$ can improve the accuracy of the estimator $\widehat\theta^K$, in the sense of reducing its uncertainty.

While one might, in principle, seek a criterion on $\Mb(\theta,\mu)$ to be optimised over all probability measures in $\P(\V)$, this generality is in fact unnecessary.
Atomic measures play a special role here: these are measures $\mu$ for which
there exists a positive integer $\ell \in \N$, inputs $(u^{(1)},\dots, u^{(\ell)})\in \V^{\ell}$, and weights $(w_1,\dots,w_{\ell})\in\R^{\ell}$, $\sum_{k=1}^{\ell} w_k =1$, such that  $\mu=\sum_{k=1}^{\ell} w_k \delta_{u^{(k)}}$.

It turns out that any Fisher information matrix achievable by an arbitrary measure can also be obtained using an atomic measure, and the value of $\ell$ has an a priori upper bound,  as stated in the following result. A proof of this fact is provided in appendix.
\begin{lemma}\label{L:conv}
Let $\mathcal{M}(\theta, \V)=\{\Mb(\theta,\mu)\mid \mu \in \P(\V)\}$ be the set of all possible Fisher matrices generated on the basis of a probability measures supported on $\V$.
Then
$$
\mathcal{M}(\theta, \V) = \Conv\{\Mb(\theta,\delta_u)\mid u\in \V\}.
$$
As a consequence, any element of $\mathcal{M}(\theta, \V)$ is the convex sum of at most $p(p+1)/2+1$ elements of the form $\Mb(\theta,\delta_u)$, $u\in \V$.
\end{lemma}

A consequence of Lemma~\ref{L:conv} is that an optimal experiment, for a given criterion to be specified, only consists of, at most, $\ell\leq p(p+1)/2+1$ different input signals $u^{(i)}$, each one of these elementary experiments being repeated a number of times proportional to the weight the optimal measure allocates to this particular $u^{(i)}$. The construction of an optimal measure is considered in the next section. 

\subsection{Optimised Fisher information}

In the present study, we focus on \textit{locally optimal design}, that is optimisation of the Fisher information matrix in order to minimise the covariance of the distribution in \eqref{E:TCL} around a nominal value $\theta^0$ that we consider a stand-in for the unknown $\bar\theta$. We consider the $D$-optimality criterion, for which an optimal experiment is obtained by the maximisation over $\mu\in \P(\V)$ of
$$
\phi(\mu)=\Phi[\Mb(\theta^0,\mu)] = \log \det \Mb(\theta^0,\mu).
$$
{\Luc We assume that $\Mb(\theta^0,\V)$ contains at least one non-singular matrix. Other design criteria could be considered as well, see, e.g., \cite{Fedorov72,Silvey80,Pukelsheim2006} for a thorough presentation of the optimal design of experiments.}
The function $\Mb\mapsto \Phi(\Mb)=\log \det(\Mb)$ is monotone on the cone of positive definite matrices, i.e. if $A,B$ and $B-A$ are all positive definite matrices ({\Luc $0\preceq A\preceq B$, in the sense of Loewner's ordering}), then $\log \det(A)\leq \log \det (B)$. {\Luc As a result, the maximal value of $\Phi$ must be reached for a matrix $\Mb$ at the boundary of $\Mb(\theta^0,\V)$.
Any element of the boundary can be written as the convex sum of $p(p+1)/2$ matrices $\Mb(\theta^0, \delta_u)$, hence the reduction by 1 of the bound of Lemma~\ref{L:conv}. The function $\Phi$ is also strictly concave, and therefore there exists a unique optimal matrix $\Mb^*$ in $\Mb(\theta^0,\V)$; note that it does not imply the existence of a {\em unique} measure $\mu^*$ such that $\Mb^*=\Mb(\theta^0,\mu^*)$.}

%\LP{On pourrait fusisonner la première phrase avec la remarque~3.2} The conclusion of  Lemma~\ref{L:conv} is that optimality can be achieved by a mixture of at most $p(p+1)/2+1$ inputs, and more are unnecessary from the $D$-optimality point of view. 
%\LP{Je ne comprends pas bien la suite du paragraphe.}
%Furthermore, an implicit consequence is that the lemma is most useful when $\V$ is not  assumed continuous, being discrete or with gaps. When $\V$ is a continuum (or contains one), the continuity of $\Mb(\theta^0, \delta_u)$ with respect to $u$ ensures that the set $\mathcal{M}(\theta^0, \V)$ is the convex hull of a continuous image, and thus the convex sum intended by Caratheodory theorem may be unhelpful (in the sense that facets of the convex may be essentially point-like). On the other hand, this fact is particularly useful when searching for optimal Fisher information matrices on a finite (but potentially large) discrete set of candidate controls, where convex analysis can be applied. This is the perspective we adopt in the design of the algorithm presented here.


%\begin{remark}
%We can reduce the a priori bound on $\ell$ by $1$. Since $\mathcal{M}(\theta^0,\V)$ is convex, the maximum of $\log\circ\det$ over $\mathcal{M}(\theta^0,\V)$ is unique and reached on the boundary $\partial \mathcal{M}(\theta^0,\V)$  (again, by monotonicity). As such, this boundary element  is actually a convex sum $p(p+1)/2$ (at the most) elements of the form $\Mb(\theta^0,\delta_u)$. By linearity of $\Mb(\theta^0,\cdot)$, the optimum of $\Phi$ is reached by an atomic measure with $\ell\leq p(p+1)/2$.
%\end{remark}


From now on, we assume that we are searching for the optimum atomic measure $\mu$, $\ell\leq p(p+1)/2$, with atoms to be selected among a finite set  $\V=\{u^{(1)}, \dots ,u^{(L)}\}$ of cardinal $|\V|=L$, $L>p(p+1)/2$.
Let
$$
\S_L=\left\{
(w_1,\dots, w_L)\in [0,1]^L
\mid 
\sum_{k=1}^L w_k=1 
\right\}
$$ 
be the 
probability
simplex in $\R^L$. %Letting $\{u^{(1)}, \dots ,u^{(L)}\}=\V$, 
We are looking for an optimal measure among the collection of measures $\mu=\sum_{k=1}^L w_k\delta_{u^{(k)}}$, $(w_1,\dots, w_L)\in \S_L$, and the goal is thus to find an optimal weight vector $w^*$.
{\Luc We know that at the optimum at least $L-p(p+1)/2$ weights $w_k^*$ can be set to zero and this will be used in Section~\ref{S:elimination} to simplify the optimisation problem, but for the moment we describe a general search method for this optimal weight distribution}.

The method relies on a multivariate extension of Kiefer-Wolfowitz equivalence Theorem; see \cite{KieferW60} for the original formulation and \cite{HarmanT2009} for the multivariate response case considered here. 
This result states that a probability measure $\mu^*=\sum_{k=1}^L w_k^*\delta_{u^{(k)}}$ is optimal, i.e.\ satisfies 
\begin{equation}\label{mu*optimal}
\phi(\mu^*)=\max_{w\in \S_L}
\left\{ \phi(\mu) \mid \mu=\sum_{k=1}^L w_k\delta_{u^{(k)}}\right\},
\end{equation}
if and only if 
\begin{equation}\label{ET}
w^*=\arg\min_{w\in \S_L}
\left\{
\max_{k=1,\dots ,L} 
\Tr
\left[
\Mb(\theta^0,\mu)^{-1}
\Mb_k
\right]
\mid 
\mu=\sum_{k=1}^L w_k\delta_{u^{(k)}}
\right\},
\end{equation}
where  $\Mb_k=\Mb(\theta^0,\delta_{u^{(k)}})$. Moreover, $\mu^*$ satisfies 
%with minimum for the above functional:
\begin{equation}\label{E:def_mu_star}
\max_{k=1,\dots ,L} 
\Tr
\left[
\Mb(\theta^0,\mu^*)^{-1}
\Mb_k
\right]=p.
\end{equation}
%
It should be understood that 
$
\Tr
\left[
\Mb^{-1}
\Mb_k
\right]- p$ is essentially the directional derivative of the strictly concave function $\log \circ \det$ at $\Mb$ in the direction $\Mb(\theta^0,\delta_{u^{(k)}})-\Mb$.
%(indeed $(\log \det M(t))' = \Tr(M(t)^{-1}M'(t))$). 
%Hence at the unique maximum of $\log\circ \det$ on the convex generated by $M_k$, all directional directional derivatives towards a point of the convex must be negative and thus  \eqref{E:def_mu_star} is true. At any other point of the convex this directional derivative is positive in at least one direction (by concavity) and $\max_{k=1,\dots, L} 
{\Luc The equivalence between \eqref{mu*optimal} and \eqref{ET} simply expresses that at
the unique maximum $\Mb^*$ of $\Phi$ over $\mathcal{M}(\theta^0, \V)$ all directional derivatives towards a matrix $\Mb_k$ must be negative, and this also implies 
\eqref{E:def_mu_star}. At any other matrix in $\mathcal{M}(\theta^0, \V)$ this directional derivative is positive in at least one direction (by concavity) and $\max_{k=1,\dots, L} 
\Tr
\left[
\Mb(\theta^0,\mu^*)^{-1}
\Mb_k
\right]>p$.  }

To find the optimal $\mu$, we use Fedorov's algorithm (\cite{Fedorov70}, \cite[Sect.~5.2]{Fedorov72}) which proposes to solve an ascent method by following the largest directional derivative. That is, at a given element $w\in \S_L$ and $\mu = \sum_{k=1}^L w_k\delta_{u^{(k)}}$, we pick
\begin{equation}\label{k*}
k^*=\arg\max_{k=1,\dots, L}  \Tr
\left[
\Mb(\theta^0,\mu)^{-1}
\Mb_k
\right],
\end{equation}
and update $w$ into $\alpha e_{k^*}+(1-\alpha)w$, 
{\Luc where $e_k$ denotes the $k$-th basis vector of $\R^L$ (with all components zero except the $k$-th one which equals one).}
%with $e_k\in \S_L$ to be $0$ except in position $k$, where it is $1$. 
{\Luc The value of $\alpha$ can be either fixed or optimised (numerically). Indeed, using the strict concavity of $\Phi$, there exists a single $\alpha^*$ that maximises $\phi[\alpha e_{k^*}+(1-\alpha)w]$ with respect to $\alpha\in[0,1]$; $\alpha^*$ can easily be determined by line search (e.g., with the Golden-Section algorithm).}  
%adaptive and (numerically) computed. 
%Indeed, using the concavity $\log \circ \det$, there exists a single value of $\alpha$ such that $\frac{\partial}{\partial \alpha}\Phi(\theta^0,\alpha \delta_{u^{(k^*)}}+(1-\alpha)\mu)=0$, at which point 
%$$
%\Tr
%\left[
%\Mb(\theta^0,\alpha \delta_{u^{(k^*)}}+(1-\alpha)\mu)^{-1}
%(\Mb_k-\Mb(\theta^0,\mu))
%\right]=0.
%$$
%This point can be, for instance, found by bisection.

\subsection{Elimination strategies}\label{S:elimination}

\paragraph{Elimination of $p(p+1)/2$-dimensional faces.}
{\Luc In the case of a structurally locally identifiable model %with a small number of parameters ($p=2,3$), 
it is expected that a single input $u$ provides enough excitation to make the parameters $\theta$ (locally) estimable, that is, most matrices $\Mb(\theta^0,\delta_{u_k})$, $k=1,\cdots,L$ are positive definite, instead of being simply semidefinite in the general case. It is then possible to further exploit concavity to accelerate the search for an optimal measure $\mu^*$.

We know that the unique optimal matrix 
$\Mb^*$
%$\Mb(\theta^0,\mu^*)$ 
belongs to the boundary of the convex set $\mathcal{M}(\theta^0,\V)$, hence it belongs to one of its  $p(p+1)/2$-dimensional faces. 
We propose a test to eliminate faces that cannot support an optimal measure $\mu^*$. Let $(M_1,\dots , M_\ell)$ be the $\ell$ matrices defining such a face, $\ell = p(p+1)/2$, with $M_k=\Mb(\theta^0,\delta_{u^{(k)}})$ for all $k$, and for 
any $w\in\S_\ell$ denote $\Mb(w;M_1,\dots , M_\ell)=\sum_{k=1}^\ell w_k M_k$.
From the concavity of $\Phi$, we have, for any full-rank matrix $\Mb\in\mathcal{M}(\theta^0,\V)$,
\begin{eqnarray*}
\log\det \left[\Mb(w;M_1,\dots, M_\ell)\right] &\leq& \log\det(\Mb) + \Tr\left[\Mb^{-1}\Mb(w;M_1,\dots, M_\ell)\right] - p \\
&\leq& \log\det(\Mb) -p  + \max_{k=1,\ldots,\ell} \Tr(\Mb^{-1}M_k) \,.
\end{eqnarray*}
This implies in particular (take $\Mb=M_i$ for some $i\in\{1,\ldots,L\}$),
\begin{eqnarray*}
\max_{w\in\S_\ell} \log\det \left[\Mb(w;M_1,\dots, M_\ell)\right] &\leq&  \min_{i=1,\ldots,L} \Psi_i(M_1,\dots , M_\ell),
%\Psi(M_1,\dots , M_\ell) \\
%&& = \min_{i=1,\ldots,L} \left\{ \log\det(M_i) - p + \max_{k=1,\ldots,\ell} \Tr\left[M_i^{-1}M_k\right] \right\},
\end{eqnarray*}
where we denote 
$$
\Psi_i(M_1,\dots , M_\ell) = \left\{\begin{array}{ll}
    \log\det(M_i) - p + \max_{k=1,\ldots,\ell} \Tr(M_i^{-1}M_k)  &  \mbox{ if }  M_i \mbox{ has full rank,}\\
    +\infty & \mbox{ otherwise.} 
\end{array}
\right.
$$
The face defined by $(M_1,\dots , M_\ell)$ can thus be eliminated as soon as 
\begin{eqnarray} \label{elimination1}
%\Psi(M_1,\dots , M_\ell)<\log\det(\widehat\Mb)
\min_{i=1,\ldots,L} \Psi_i(M_1,\dots , M_\ell)<\log\det(\widehat\Mb)
\end{eqnarray}
for some matrix $\widehat\Mb\in\mathcal{M}(\theta^0,\V)$ (note that it does not mean that one or several $\delta_{u^{(k)}}$ that contribute to this face will not be used is another face supporting the optimal $\mu^*$). 
The matrices $\Mb_i$, $i=1,\ldots,L$ are obvious candidates for $\widehat\Mb$, and we can thus safely eliminate the face $(M_1,\dots , M_\ell)$ when
\begin{eqnarray} \label{elimination1b}
%\Psi(M_1,\dots , M_\ell) < \max_{u\in\V} \phi(\delta_u) \,.
\min_{i=1,\ldots,L} \left\{\Phi(M_i)-p+\max_{k=1,\ldots,\ell} \Delta_{i,k}\right\} < \max_{i=1,\ldots,L} \Phi(M_i) \,,
\end{eqnarray}
where the minimum in \eqref{elimination1b} is over matrices $M_i$ having full rank, and where we have denoted 
\begin{eqnarray} \label{Delta}
\Delta_{i,k}=\Tr(M_i^{-1}M_k), \ i,k=1,\ldots,L.
\end{eqnarray}

The strategy \eqref{elimination1b} can be applied independently of the optimisation algorithm, as it only uses the information provided by the input candidates $u^{(1)},\ldots,u^{(L)}$. However, the elimination of useless faces can also be considered during the search for an optimal measure. In that case, if $\mu$ is the current measure at a given iteration of the optimisation algorihm, we use $\widehat\Mb=\Mb(\theta^0,\mu)$ in \eqref{elimination1}, which provides a more efficient elimination than with the elementary matrices $M_i$. 

\vspace{0.3cm}
Even for moderate $L$, due to combinatorial explosion when considering all 
$\binom{L}{\ell}$ faces with $\ell=p(p+1)/2$ matrices,
%$p(p+1)/2$ faces for large $p$, 
the method above is restricted to small values of $p$ ($p=2,3$). In contrast, the elimination rule presented below allows us to eliminate inputs $u^{(k)}$ that cannot support a $D$-optimal experiment, using the information provided by an arbitrary measure $\mu=\sum_{k=1}^L w_k\delta_{u^{(k)}}$, i.e., an arbitrary vector $(w_1,\dots, w_L)\in \S_L$, associated with a non-singular matrix $\Mb\in\mathcal{M}(\theta^0,\V)$. 
}

\begin{comment}
We propose a test to determine which facets may be admissible to carry the maximum of $\Phi$. For this we use a test using only the values at the vertices of the convex. Let $M_k=\Mb(\theta^0,\delta_{u^{(k)}})$, $k=1,\dots,L$. Since $\Phi$ is concave, we know that for any $w\in \S_L$
\begin{equation}
    \sum_{k=1}^L w_k \log\det(M_k)\leq \log\det\left(\sum_{k=1}^L w_k M_k\right)
    \leq 
    \Phi^*=\max_{w\in \S_L}
\log\det\left(\sum_{k=1}^L w_k M_k\right).
\end{equation}
In particular,
\begin{equation}
    \max_{k\in \{1,\dots,L\}} \log\det(M_k)\leq 
    \Phi^*.
\end{equation}
On the other hand, we can also use the concavity of $\Phi$ to get an upper bound on $\Phi^*$.
Indeed, for  for any $l\in \{1,\dots,L\}$, for any $N\in \R^{p\times p}$, $\dfrac{\diff}{\diff t}_{|t=0}\log \det (M_l+ t N)=\Tr\left(M_l^{-1} N \right)$ (assuming, again, definite-positiveness).

Since concave functions are below their tangent lines, using a sub-collection $(M_1,\dots, M_\ell)$ and
applying this fact with $N=\left(\sum_{k=1}^\ell w_k M_k\right)-M_l$
yields
\begin{equation}
\log\det\left(\sum_{k=1}^\ell w_k M_k\right)
\leq 
\log\det\left(M_l\right)
+
\Tr\left(M_l^{-1}
\left(\sum_{k=1}^K w_k M_k\right)
\right)-p.
\end{equation}
Hence 
\begin{equation}
\begin{aligned}
\log\det\left(\sum_{k=1}^\ell w_k M_k\right)
&\leq 
\min_{l\in \{1,\dots,\ell\}}
\left[
\log\det\left(M_l\right)
-p
+
\sum_{k=1}^\ell w_k
\Tr\left(
M_l^{-1}
 M_k
\right)
\right]
\\
&\leq 
\min_{l\in \{1,\dots,\ell\}}
\left[
\log\det\left(M_l\right)
-p
+
\max_{k\in \{1,\dots,\ell\}}
\Tr\left(M_l^{-1}
 M_k\right)
\right].
\end{aligned}
\end{equation}



Now if we consider a given facet of $\mathcal{M}(\theta^0,\V)$ with vertices $(M_1,\dots , M_\ell)$, $\ell = p(p+1)/2$, we can give a priori bounds on the optimal value of $\Phi$ within the facet. Letting 
\begin{equation}
\Psi(M_1,\dots , M_\ell)=\min_{l\in \{1,\dots,K\}}
\left[
\log\det\left(M_l\right)
-p
+
\max_{k\in \{1,\dots,\ell\}}
\Tr\left(M_l^{-1}
 M_k\right)
\right],
\end{equation}
it is clear that 
\begin{equation}
    \max_{k\in \{1,\dots,K\}} \log\det(M_k) \leq\Phi^* \leq \Psi(M_1,\dots , M_\ell).
\end{equation}
This inequality has an important consequence. If
\begin{equation}
    \Psi(M_1,\dots , M_\ell)<\max_{u\in \V} \Phi(\delta_u),
\end{equation}
then the facet with vertices $(M_1,\dots , M_\ell)$ does not contain the maximum of $\log\det$ over the convex $\mathcal{M}(\theta^0, \V)$. 
\end{comment}

\paragraph{Elimination of elementary matrices $M_i$.} 
{\Luc 
%Due to combinatorial explosion when considering all $p(p+1)/2$ faces for large $p$, the method above is restricted to small values of $p$ ($p=2,3$). In contrast, the elimination rule of \cite{HarmanT2009} allows us to eliminate inputs $u^{(k)}$ that cannot support a $D$-optimal experiment, using the information provided by an arbitrary measure $\mu=\sum_{k=1}^L w_k\delta_{u^{(k)}}$, i.e., an arbitrary vector $(w_1,\dots, w_L)\in \S_L$, associated with a non-singular matrix $\Mb\in\mathcal{M}(\theta^0,\V)$. 

Let $\Mb$ be a non-singular matrix in $\mathcal{M}(\theta^0,\V)$.
Denote $\epsilon(\Mb)=\Tr( \Mb^{-1}M_{k^*} )-p$, where $k^*=k^*(\Mb)$ is given by \eqref{k*} and $M_k=\Mb(\theta^0,\delta_{u^{(k)}})$, $k=1,\ldots,L$. From \cite[Th.~3]{HarmanT2009}, any $u^{(k)}$ such that
\begin{equation}\label{elimination2}
    \Tr(\Mb^{-1}M_k) < p \left[1+\frac{\epsilon(\Mb)}{2} - \frac{\sqrt{\epsilon(\Mb)[4+\epsilon(\Mb)-4/p]}}{2} \right]
\end{equation}
can be safely eliminated from the set $\V$ of input candidates; see also \cite{HPa06_SPL}. The elimination becomes more an more efficient as $\Mb$ approaches $\Mb^*$, hence the interest of coupling this elimination strategy with an optimisation algorithm. 

The same approach can also be used prior to optimisation, by susbtituting the matrices $M_i$ in turn, $i=1,\ldots,L$, for $\Mb$. Denote $\epsilon_i=\max_{k=1,\ldots,L} \Delta_{i,k}-p$ with $\Delta_{i,k}$ given by \eqref{Delta}. Then, \eqref{elimination2} implies that any $u^{(k)}$ such that 
$$
\min_{i=1,\ldots,L} \left\{ \Delta_{i,k} - p \left[1+\frac{\epsilon_i}{2} - \frac{\sqrt{\epsilon_i(4+\epsilon_i-4/p)}}{2} \right] \right\} <0
$$
can be safely eliminated as it cannot support an optimal measure $\mu^*$. As in \eqref{elimination1b}, the minimum is with respect to matrices $M_i$ of rank $p$ (when $M_i$ approaches singularity, $\epsilon_i$ tends to $+\infty$, $p [1+\epsilon_i/2 - (1/2)\sqrt{\epsilon_i(4+\epsilon_i-4/p)}]$ tends to 1, and a $M_i$ close to singularity only eliminates matrices $M_k$ such that $\Tr(M_i^{-1}M_k) \lesssim 1$). 


\LP{Je crois que ça vaudrait le coup d'essayer ça et de comparer avec l'élimination des faces complètes du paragraphe précédent (il n'y a que $L^2$ calculs de $\Delta_{i,k}$ à effectuer).}
}

\subsection{Computing the Fisher information matrix}\label{S:computing-FIM}

The crucial observation is that if $x$ is a solution of $\dot x=A(\theta) x+Bu$, then (recall we assumed $\theta\mapsto A(\theta)$ to be of regularity $C^1$ at least), for all $1\leq i\leq p$,
$$
\frac{\diff}{\diff t} \partial_{\theta_i} x= \frac{\partial}{\partial \theta_i} \frac{\diff x}{\diff t}= \partial_{\theta_i} A(\theta)x + A(\theta) \partial_{\theta_i} x.
$$
We may write as well
$$
\frac{\diff}{\diff t}
\begin{pmatrix}
x
\\
\partial_{\theta_i} x
\end{pmatrix}
=
\bar A_i
\begin{pmatrix}
x
\\
\partial_{\theta_i} x
\end{pmatrix}
+
\bar B u,
$$
%From this observation, following steps similar to that of section~\ref{S:dynamics}, the construction of the derivatives of each element with respect to the parameters follows.
letting 
$\bar A_i(\theta)=\begin{pmatrix}
    A(\theta) & 0
    \\
    \partial_{\theta_i}A(\theta) & A(\theta)
\end{pmatrix}$,
$\bar B=\begin{pmatrix}
    B
    \\
    0
    \end{pmatrix}$ and, for later, $\bar C=\begin{pmatrix}
0&C
\end{pmatrix}$.
This observation still applies for SDE systems of the type of \eqref{E:sde1}.
Let $\bar W$ be a standard $2n$-dimensional Brownian motion. With $\bar x^i=\begin{pmatrix}
x
\\
\partial_{\theta_i} x
\end{pmatrix}$,
the controlled stochastic differential equation becomes:
\begin{equation}\label{E:sde2}
\diff \bar x^i(t)
=
\bar A_i(\theta)
\bar x^i(t)
+
\bar B u(t)
+
\begin{pmatrix}
D & 0
\\
0 & 0
\end{pmatrix}\diff \bar W(t)
\end{equation}
with initial condition $(x_0,0)$.
As before, increments of $\bar x^i$ between two dates $t_{k},t_{k+1}$, $0\leq k<N$, satisfy
\begin{equation}\label{E:rec1}
\bar x^i_{k+1}=\bar F_i \bar x^i_k
+
\bar G_i u_k
+
\bar W^i_k
\end{equation}
with
\begin{equation}
\bar F_i=e^{\Delta \bar A_i(\theta)}
\qquad 
\bar G_i
= 
\int_{0}^{\Delta}e^{\bar A_i(\theta)(\Delta-s)}\bar BU(s)\diff s,
\quad
\bar W_k^i=\int_{t_k}^{t_{k+1}}e^{\bar A_i(\theta)(t_{k+1}-s)}\begin{pmatrix}
D & 0
\\
0 & 0
\end{pmatrix} \diff \bar W(s).
\end{equation}
As before $(\bar W_k^i)_{0\leq k<N}$ is an i.i.d. sequence of centered Gaussian random vectors with covariance matrix $\bar\Sigma^i$, i.e. $\bar W_k^i\sim \mathcal{N}(0,\bar\Sigma^i)$, with:
$$
\bar\Sigma^i = \int_0^\Delta e^{\bar A_i(\theta)(\Delta-s)}
\begin{pmatrix}
D D^\top  & 0
\\
0 & 0
\end{pmatrix} 
e^{\bar A_i^\top(\theta)(\Delta-s)}\diff s.
$$
The conclusion is that 
\begin{equation}\label{eq:deta_dtheta}
\partial_{\theta_i}\eta(t_k,\theta,u)
=
\bar C
{\bar F_i}^{k}
\begin{pmatrix}
x_0\\0
\end{pmatrix}
+
\sum_{l=0}^{k-1} 
\bar C
{\bar F_i}^{k-l-1}{\bar G_i}u_{l} 
\end{equation}
and
\begin{equation}\label{eq:dy_dtheta}
\partial_{\theta_i}y_k(u)
=\partial_{\theta_i}\eta(t_k,\theta,u)+\sum_{l=0}^{k-1} \bar C {\bar F_i}^{k-l-1}{\bar W^i_l}.
\end{equation}

In addition to being able to differentiate the output, this construction also allows to differentiate $\Sigmab$. Indeed, looking at $\Cov(y_k,y_l)$, we have
\begin{multline}
\partial_{\theta_i }
\Cov(y_k,y_l)
=
\Cov(\partial_{\theta_i }y_k,y_l)
+
\Cov(y_k,\partial_{\theta_i}y_l)
=
\\
\begin{pmatrix}
0 & C
\end{pmatrix}
\Cov(\bar x^i_k, \bar x^i_l)
\begin{pmatrix}
C^\top \\ 0
\end{pmatrix}
+
\begin{pmatrix}
C & 0
\end{pmatrix}
\Cov(\bar x^i_k, \bar x^i_l)
\begin{pmatrix}
0 \\ C^\top 
\end{pmatrix}.  
\end{multline}
Hence, with $\bar C=\begin{pmatrix}
0&C
\end{pmatrix}$ as before and $\bar C'=\begin{pmatrix}
C & 0
\end{pmatrix}$, we define $\bar \Sigma_l^i=\sum_{j=0}^{l-1}\bar F_i^{j}\bar \Sigma^i(\bar F_i^j)^{\top}$ and 
$$
\bar \Sb^i
=
\left(\Cov(\bar x^{i}_k, \bar x^{i}_l)\right)_{1\leq k,l\leq N}
=
\begin{pmatrix}
\bar \Sigma_1^i & \bar \Sigma_1^i \bar F_i^\top & \bar \Sigma_1^i (\bar F^2_i)^\top & \cdots 
\\
\bar F_i \bar \Sigma_1^i & \bar \Sigma_2^i  & \bar \Sigma_2^i \bar F^\top_i & \cdots 
\\
\bar F^2_i \bar \Sigma_1^i & \bar F_i \bar \Sigma_2^i  & \bar \Sigma_3^i  & \cdots 
\\
\vdots&\vdots&\vdots&\ddots
\end{pmatrix},
$$
\begin{equation}
\frac{\partial \Sigmab}{\partial \theta_i}
=
\diag(\bar C,\dots, \bar C)
\cdot \Sb \cdot
\diag(\bar C',\dots, \bar C')^\top
+
\diag(\bar C',\dots, \bar C')
\cdot \Sb \cdot
\diag(\bar C,\dots, \bar C)^\top.
\end{equation}
This allows to compute the matrix $Q$ from \eqref{E:defM}-\eqref{E:defQ}.



% \begin{lemma}
% Let $B'=\begin{pmatrix}
%     B
%     \\
%     0
%     \end{pmatrix}$ and $C'=\begin{pmatrix}
% 0&C
% \end{pmatrix}$.
% For $i=1,\dots, p$, let  
% $A'_i(\theta)=\begin{pmatrix}
%     A(\theta) & 0
%     \\
%     \partial_{\theta_i}A(\theta) & A(\theta)
% \end{pmatrix}$,
% $F'_i=\exp[ \Delta A'_i(\theta)]$
% and
% \begin{equation}
%     G'_i= 
%     \int_{0}^{\Delta}
%     \exp[ (\Delta-s)A'_i(\theta)]B'
%     U(s)\diff s.
% \end{equation}


% Then
% \begin{equation}\label{eq:deta_dtheta}
% \partial_{\theta_i}\eta(t_k,\theta,u)
% =
% C'
% {F'_i}^{k}
% \begin{pmatrix}
% x_0\\0
% \end{pmatrix}
% +
% \sum_{l=0}^{k-1} 
% C'
% {F'_i}^{k-l-1}{G'_i}u_{l} .
% \end{equation}
% \end{lemma}


Computationally, we have
\begin{equation}
    \bar{F}_i=
    \begin{pmatrix}
    F & 0
    \\
    \partial_{\theta_i} F &  F
    \end{pmatrix}.
\end{equation}
As a consequence we have
\begin{equation}
\bar F_i^k
=
\begin{pmatrix}
    F^k & 0
    \\
    \sum_{l=0}^k F^l (\partial_{\theta_i} F) F^{k-l} &  F^k
    \end{pmatrix},
\end{equation}
so only powers of $F$ are necessary to compute powers of $\bar F_i$.
Set 
\begin{equation}
H_i
=
    \diag(\bar C,\dots, \bar C)
    \begin{pmatrix}
    \id_{\R^{2n\times 2n}} &  &  &
    \\
    \bar F_i & \ddots &  &(0)
    \\
    {\bar F_i}^2  &\ddots  &\ddots& &
    \\
    \vdots &\ddots  &\ddots&\ddots &
    \\
    {\bar F_i}^{N-1}&\cdots&{\bar F_i}^2&{\bar F_i}&\id
    \end{pmatrix}
    \diag(\bar G_i,\dots, \bar G_i) \in \R^{Nq\times N m},
\end{equation}
and
\begin{equation}
    K_i
=
    \begin{pmatrix}
    \bar C \bar F_i 
    \\
    \bar C \bar F_i^2
    \\
    \vdots 
    \\
   \bar C{\bar F_i}^{N}
    \end{pmatrix}
    \begin{pmatrix}
    \id
    \\
    0_{\R^{n\times n}}
    \end{pmatrix}
    \in \R^{Nq\times n}
\end{equation}
Denoting $\ub\in \R^{Nm}$ the vector concatenation of $u_0,\dots ,u_{N-1}$, that is, $\ub^\top=\begin{pmatrix}
u_0^\top & \cdots & u_{N-1}^\top
\end{pmatrix}$,
\begin{equation}
\partial_{\theta_i}\etab(\theta,u)=H_i \ub + K_i x_0.
\end{equation}
Then we get the full expression of the Fisher information matrix $\Mb$:
\begin{equation}
    \Mb_{ij}(\theta,\delta_{u})
    =
    \left(H_i \ub + K_i x_0\right)^\top \Sigmab^{-1}\left(H_j \ub + K_j x_0\right) + Q_{ij}.
\end{equation}
In particular setting $L_{ij}=(H_i^\top \Sigmab^{-1} H_j+H_j^\top \Sigmab^{-1} H_i)/2$ and $x_0=0$,
%$L'_{ij}=(H_i^\top \Sigmab^{-1} K_j+H_j^\top \Sigmab^{-1} K_i)/2$ and
%$L''_{ij}=(K_i^\top \Sigmab^{-1} K_j+K_j^\top \Sigmab^{-1} K_i)/2$, 
we finally get
\begin{equation}
    \Mb_{ij}(\theta,\delta_{u})
    =
    \ub^\top L_{ij} \ub+ Q_{ij}.
\end{equation}


\begin{remark}\label{R:approximate-design}
Note that in this case the implementation of the whole experiment defined by the measure $\mu=\sum_{k=1}^K w_k\delta_{u^{(k)}}$ only requires rescaling each $u^{(k)}$ by the factor $\sqrt{Kw_k}$ in the $k$-th elementary experiment with $u^{(k)}$; that is,
$$
\Mb(\theta,\mu) = \frac1K\, \sum_{k=1}^K \Mb(\theta,\delta_{\sqrt{Kw_k}u^{(k)}}) \,.
$$
%The set $\V$ used for the construction of an optimal measure $\mu$ can be such that all $u\in\V$ satisfy the power normalisation $\ub^\top\ub=1$, and the 
The asymptotic results of Section~\ref{S:FIM} correspond to the situation where the number of repetitions of the whole experiment associated with $\mu$, involving $K$ different input signals, tends to infinity.
%
The application of approximate design theory, that is, the definition of an experiment as a probability (design) measure, here does not imply any approximation or rounding when implementating the experiment, in contrast with \cite{PukelsheimR92}. This situation is more commonly encountered when designing an optimal input signal in the frequency domain for the identification of a linear dynamical system; see \cite[Chap.~6]{GoodwinP77}, \cite{Zarrop79}.
\end{remark}


\section{Least-squares estimation}

We now discuss the optimisation step for estimation of the parameter $\theta\in \Theta\subset \R^p$. Once $\mu^*$ has been selected, the goal is to find an estimate of $\theta$ by solving 
\begin{equation}
    \hat \theta=\arg\max_{\theta\in \Theta}\Jb(\theta,\mu^*).
\end{equation}
We propose to solve this method by gradient descent.

\subsection{Direction of descent}
We have
\begin{equation}
\Jb(\theta,\mu)
=
\int_\V
(\yb(u)-\etab(\theta,u))^\top 
\Sigmab^{-1}(\theta)
(\yb(u)-\etab(\theta,u))\diff \mu(u)+\log\det \Sigmab(\theta).
\end{equation}
Hence the derivative of $\Jb(\theta,\mu)$ with respect to $\theta_i$ is given by:
\begin{equation}\label{E:dJdtheta_i}
\begin{aligned}
\frac{\partial \Jb}{\partial\theta_i}(\theta,\mu)
=&
-2\int_\V
(\yb(u)-\etab(\theta,u))^\top 
\Sigmab^{-1}(\theta)
\frac{\partial\etab(\theta,u)}{\partial \theta_i}\diff \mu(u)
\\
&+
\int_\V
(\yb(u)-\etab(\theta,u))^\top 
\Sigmab^{-1}(\theta) \frac{\partial \Sigmab}{\partial\theta_i}\Sigmab^{-1}(\theta)
(\yb(u)-\etab(\theta,u))\diff \mu(u)
\\&
+
\Tr\left(\Sigmab^{-1}(\theta) \frac{\partial \Sigmab}{\partial\theta_i}\right).
\end{aligned}
\end{equation}

\section{Linear model of the retina}

\subsection{Context}\label{S:Context}
 
The general context of this paper concerns dynamical systems of the form \eqref{E:sde1} below, that is $n$-dimensional and non autonomous linear stochastic equations. The linear part is controlled by a $n \times n$ matrix, $A$, depending on a set of $p$ parameters i.e.\ a $p$ dimensional vector $\theta$.
The non autonomous part depends on a stimulus (or control) $u(t)$ modulated by a matrix term $B$. The stochastic part is a Brownian noise tuned by a covariance matrix $D$. As explained in the introduction, the equation   \eqref{E:sde1} is motivated by recent works attempting to model the retinal network and its response to general stimuli \cite{souihel-cessac:21,cessac:22,kartsaki-hilgen-etal:24,ebert-buffet-etal:24,emonet-souihel-etal:25}. In these papers, the matrix $A$ has a specific structure that mimics the connectivity between retinal layers. In short, the retina is structured into successive layers of neurons processing visual scenes: photoreceptors (PRs) which converts photons into variations of electric potentials; bipolar cells (BCs) that receive the input from photoreceptors and are modulated by horizontal cells (HCs); amacrine cells (ACs) that receive and modulate the input from bipolar cells and modulate retinal ganglion cells (RGCs), the last layer in the retina. The axons of ganglion cells constitute the optic nerve which sends the retina output to the brain in the form of spike trains. In many models, the conjugated effect of PRs, BCs and HCs is modelled by a convolution kernel followed by a non linearity \B{Add refs}. In general, many nonlinear effects take place in the retina but, here, we are going to neglect them, for simplicity.

Typically, in the model, the structure of $A$ is guided by the physiology of the retina. 
Practically, when processing experimental data, the dimension of $A$ is from a few hundreds to thousands. The parameters vector $\theta$ encompasses characteristic integration times of cell types and the amplitude of the synaptic strengths. Virtually the dimension $p$ of this vector can be huge, but, as shown in \cite{kartsaki-hilgen-etal:24} it is possible to reduce its dimensionality by considering averaged quantities so that, at least as a starting point, $p$ can be reduced to about $10$.  The stimulus $u$ has, in most experiments, a specific structure, coming from the physiology: it only inputs BCs (and the matrix B can be viewed as a modelling of the convolution that takes place between PRs, HCs and BCs). Thus, many of its entries are zero, although the formalism developed below does not stick at this specificity. Finally, the matrix $D$ shaping the noise term mimics correlations observed in spontaneous activity (without stimulus) {\Luc which are called noise correlations in the neuroscience literature}. 
%and called in the neuroscience literature, noise correlations. 
 
This given, the objectives, as sketched in the introduction, are the following:
 %
 \begin{enumerate}
     \item Find an efficient algorithmic method to estimate the parameter $\theta$ from experimental observations, $B$ and $D$ being given, in a large dimensional problem where $n \sim 1000$.\LP{Je crois qu'il faudrait parler des observations et de la matrice $C$.}
     \item Use optimal design theory to develop an efficient algorithm allowing us to determine an optimal stimulus in ``real time".
     \item Here we are considering an optimality criterion based on Fisher matrix. An interesting question, addressed in the conclusion, is how this criterion can be interpreted in a neuroscience context ?
 \end{enumerate}


\section{Simulations}

Pick $A,B,C,D$ as such.

Let $h\in \N$ and 
$$
\Gamma = 
\begin{pmatrix}
0 & 1 &(0)
\\
1 & \ddots & \ddots
\\
(0)& \ddots&\ddots
\end{pmatrix}\in \R^{h\times h}.
$$
We set $n=3h$ and $\theta = (\delta_1,\delta_2,\delta_3,\alpha,\beta,\gamma_1,\gamma_2)\in \R^7$, so that (with $\id:=\id_{\R^{k}})$
$$
A=
\begin{pmatrix}
-\delta_1 \id & \alpha \Gamma & 0_{\R^{h\times h}}
\\
\beta \Gamma & -\delta_2 \id & 0_{\R^{h\times h}}
\\
\gamma_1 \id & \gamma_2 \id & -\delta_3 \id
\end{pmatrix}
$$
Set $m=q=1$, and with $b\in \R^{h\times 1}$, $c\in \R^{1\times h}$, $d\in \R$
$$
B=
\begin{pmatrix}
b
\\
0_{\R^{2h\times 1}}
\end{pmatrix},\qquad
C
=
\begin{pmatrix}
c
&
0_{\R^{1\times 2h}}
\end{pmatrix},
\qquad D= d \; \id_{\R^n}.
$$



\bibliographystyle{abbrv}
\bibliography{biblio}


\appendix
\section{Appendices}
\subsection{Proofs}

For completeness, we first restate a fairly immediate fact. 
\begin{lemma}
With $\V$ compact, for any $\mu\in \P(\V)$, there exists a sequence $(\mu_n)$ in $\P(\V)$ such that each $\mu_n$  is supported on a finite number of points in $\V$ and $\mu_n\xrightarrow[n\to\infty]{w*}\mu$.
\end{lemma}
\begin{pf}
Let $\mu\in \P(\V)$.
For any positive integer $n$, we define the open cover of $\V$ given by $\V\subset \cup_{u\in V} B(u,1/n)$ (where $B(u,1/n)$ denotes the open ball of center $u$ and radius $\frac{1}{n}$). By compactness of $\V$, there exists a finite collection $u_1,\dots, u_{k(n)}\in \V$ such that $\V\subset \cup_{i=1}^{k(n)} B(u_i,1/n)$. Then we set $\mu_n=\sum_{i=1}^{k(n)} \mu(\V\cap B(u_i,1/n))\delta_{u_i}$. Now we can prove that $\mu_n\xrightarrow[n\to \infty]{w*}\mu$.

Let $f\in C^0(\V,\R)$. Since $\V$ is compact, $f$ is uniformly continuous and for all $\varepsilon>0$, there exists $\delta>0$ such that $|f(x)-f(y)|<\varepsilon$ as soon as $|x-y|<\delta$. Now take $n>1/\delta$, we have
$$
\begin{aligned}
\left|\int_\V f(u)\diff \mu(u)-\int_\V f(u)\diff\mu_n(u)\right|
&\leq
\sum_{i=1}^{k(n)}
\int_{\V\cap B(u_i,1/n)} |f(u)-f(u_i)|\diff\mu(u)
\\
&\leq 
\sum_{i=1}^{k(n)}\varepsilon\mu(\V\cap B(u_i,1/n))
=\varepsilon.
\end{aligned}
$$
Hence $\mu_n\xrightarrow[n\to \infty]{w*}\mu$.
\end{pf}


\begin{pf}[Proof of Lemma~\ref{L:conv}]
Since $\P(\V)$ is compact with respect to the weak-$*$ topology against continuous functions, for any $\mu\in \P(\V)$, there exists a sequence $(\mu_n)$ in $\P(\V)$ such that each $\mu_n$  is supported on a finite number of points in $\V$ and $\mu_n\xrightarrow[n\to\infty]{w*}\mu$. The finite support of  $\mu_n$ is equivalent to $\mu_n\in \Conv \{\delta_u\mid u\in \V\}$. Now it should be noted that for a collection  of probability measures $\mu^1,\dots, \mu^K$ and scalars $0\leq \lambda_1,\dots ,\lambda_K\leq 1$, $\sum_{k=1}^K\lambda_k=1$, we have 
$$
\Mb\left(\theta,\sum_{k=0}^K \lambda_k\mu^k\right)=\sum_{k=0}^K \lambda_k \Mb\left(\theta,\mu^k\right).
$$
Thus, since $\mu_n\in \Conv \{\delta_u\mid u\in \V\}$, $\Mb(\mu_n)\in\Conv \{\Mb(\theta,\delta_u)\mid u\in \V\}$. Then the weak-$*$ convergence of $\mu_n$ to $\mu$ and continuity of $u\mapsto \frac{\partial\etab^\top(\theta,u)}{\partial\theta}
\Sigmab^{-1} \frac{\partial\etab(\theta,u)}{\partial\theta^\top}$ imply that 
$$
\Mb(\theta,\mu)\in \overline{\Conv}\{\Mb(\theta,\delta_u)\mid u\in \V\}.
$$ 

Since $\V$ is compact and $u\mapsto \Mb(\theta,\delta_u)$ is continuous, the set  $\{\Mb(\theta,\delta_u)\mid u\in \V\}$ is also compact. As a consequence, the convex hull $\Conv \{\Mb(\theta,\delta_u)\mid u\in \V\}$ is compact and, in particular, closed. Finally, $\Mb(\mu)\in \Conv\{\Mb(\theta,\delta_u)\mid u\in \V\}$.

In conclusion, applying Carathéodory's theorem, since symmetric $p\times p$ matrices form a $p(p+1)/2$ space, we get the statement.
\end{pf}




\subsection{Second-order derivative of the estimation criterion}
The techniques developed in the paper allow to compute a second variation of the functional at any $\theta$, this can be useful for a Newton method. We give the construction for completeness. First the dynamics.

We follow the same initial observation on the derivative of a trajectory with respect to a parameter. For all $1\leq i,j\leq p$,
$$
\frac{\diff}{\diff t} \frac{\partial^2 x}{\partial \theta_i\partial \theta_j}
= 
\frac{\partial^2 }{\partial \theta_i\partial \theta_j} \frac{\diff x}{\diff t}
= 
\frac{\partial^2  A(\theta)}{\partial \theta_i\partial \theta_j}x +
\frac{\partial  A(\theta)}{\partial \theta_i}\frac{\partial  x}{\partial \theta_j} +
\frac{\partial  A(\theta)}{\partial \theta_j}\frac{\partial  x}{\partial \theta_i} +
A(\theta)   \frac{\partial^2 x}{\partial \theta_i\partial \theta_j}.
$$
We may write as well
$$
\frac{\diff}{\diff t}
\begin{pmatrix}
x
\\
\partial_{\theta_i} x
\\
\partial_{\theta_j} x
\\
\partial_{\theta_i\theta_j} x
\end{pmatrix}
=
\tilde A_{ij}
\begin{pmatrix}
x
\\
\partial_{\theta_i} x
\\
\partial_{\theta_j} x
\\
\partial_{\theta_i\theta_j} x
\end{pmatrix}
+
\tilde B u,
$$
%From this observation, following steps similar to that of section~\ref{S:dynamics}, the construction of the derivatives of each element with respect to the parameters follows.
letting 
$$
\tilde A_{ij}(\theta)=\begin{pmatrix}
    A(\theta) & 0 & 0 & 0
    \\
    \partial_{\theta_i}A(\theta) & A(\theta)& 0 & 0
    \\
    \partial_{\theta_j}A(\theta) & 0 & A(\theta) & 0
    \\
    \partial_{\theta_i\theta_j}A(\theta) &\partial_{\theta_j}A(\theta) & \partial_{\theta_i}A(\theta) & A(\theta)
\end{pmatrix}\in \R^{4n\times 4n},
$$
$\tilde B=\begin{pmatrix}
    B
    \\
    0_{\R^{3n\times m}}
    \end{pmatrix}\in \R^{4n\times m}$ and, for later, $\tilde C=\begin{pmatrix}
0_{\R^{q\times 3n}}&C
\end{pmatrix}\in \R^{q\times 4n}$.

Construction \eqref{E:sde2}-\eqref{eq:dy_dtheta} still hold in this form to get second variation of the output. In particular we denote by $\tilde x^{ij}_k$ the second variation with respect to $\theta_i,\theta_j$ of the state at date $t_k$, and 
$$
\tilde\Sigma^{ij} = \int_0^\Delta e^{\tilde A_{ij}(\theta)(\Delta-s)}
\begin{pmatrix}
D D^\top  & 0_{\R^{n\times 3n}}
\\
 0_{\R^{3n\times n}} & 0_{\R^{3n\times 3n}}
\end{pmatrix} 
e^{\tilde A_{ij}^\top(\theta)(\Delta-s)}\diff s\in \R^{4n\times 4n}.
$$

The same observation as before allows to get second variation of $\Sigmab$.
Indeed, letting $\tilde C'=\begin{pmatrix}
C & 0_{\R^{q\times 3n}}
\end{pmatrix}$,
$\tilde C'_i=\begin{pmatrix}
0_{\R^{q\times n}} & C & 0_{\R^{q\times 2n}}
\end{pmatrix}$, $\tilde C'_j=\begin{pmatrix}
0_{\R^{q\times 2 n}} & C & 0_{\R^{q\times n}}
\end{pmatrix}$, we have
\begin{equation}
\begin{aligned}
\frac{\partial^2}{\partial \theta_i \partial \theta_j}
\Cov(y_k,y_l)
& =
\Cov(\partial_{\theta_i \theta_j}y_k,y_l)
+
\Cov(\partial_{\theta_i }y_k,\partial_{\theta_j}y_l)
+
\Cov(\partial_{\theta_j}y_k,\partial_{\theta_i}y_l)
+
\Cov(y_k,\partial_{\theta_i\theta_j}y_l)
\\
&=
\tilde C
\Cov(\tilde x^{ij}_k, \tilde x^{ij}_l)
\tilde C'
+
\tilde C'_i
\Cov(\tilde x^{ij}_k, \tilde x^{ij}_l)
\tilde C'_j
\\
&\quad +
\tilde C'_j
\Cov(\tilde x^{ij}_k, \tilde x^{ij}_l)
\tilde C'_i
+
\tilde C'
\Cov(\tilde x^{ij}_k, \tilde x^{ij}_l)
\tilde C.
\end{aligned}   
\end{equation}
Denoting  $\tilde \Sigma_l^{ij}=\sum_{k=0}^{l-1}\tilde F_{ij}^{k}\tilde \Sigma^{ij}(\bar \tilde F_{ij}^k)^{\top}$ and 
$$
\tilde 
\Sb^{ij}=\left(\Cov(\tilde x^{ij}_k, \tilde x^{ij}_l)\right)_{1\leq k,l\leq N}=
\begin{pmatrix}
\tilde \Sigma_1^{ij} & \tilde \Sigma_1^{ij} \tilde F_{ij}^\top & \tilde \Sigma_1^{ij} (\tilde F^2_{ij})^\top & \cdots 
\\
\tilde F_{ij} \tilde \Sigma_1^{ij} & \tilde \Sigma_2^{ij}  & \tilde \Sigma_2^{ij} \tilde F^\top_{ij} & \cdots 
\\
\tilde F^2_{ij} \tilde \Sigma_1^{ij} & \tilde F_{ij} \tilde \Sigma_2^{ij}  & \tilde \Sigma_3^{ij}  & \cdots 
\\
\vdots&\vdots&\vdots&\ddots
\end{pmatrix},
$$
we get again
\begin{equation}
\begin{aligned}
\frac{\partial^2 \Sigmab}{\partial \theta_i\partial \theta_j}
= \;&
\diag(\tilde C,\dots, \tilde C)
\cdot\Sb^{ij}\cdot
\diag(\tilde C',\dots, \tilde C')^\top
\\&+
\diag(\tilde C'_i,\dots, \tilde C'_i)
\cdot\Sb^{ij}\cdot
\diag(\tilde C'_j,\dots, \tilde C'_j)^\top
\\
&+
\diag(\tilde C'_j,\dots, \tilde C'_j)
\cdot\Sb^{ij}\cdot
\diag(\tilde C'_i,\dots, \tilde C'_i)^\top
\\&+
\diag(\tilde C',\dots, \tilde C')
\cdot\Sb^{ij}\cdot
\diag(\tilde C,\dots, \tilde C)^\top
\in \R^{Nq\times Nq}.
\end{aligned}
\end{equation}

Now we can differentiate \eqref{E:dJdtheta_i} with respect to $\theta_j$ to obtain the second variation of the loss function.
\begin{equation}
\begin{aligned}
\frac{\partial^2 \Jb}{\partial\theta_i\partial\theta_j}(\theta,\mu)
=&\;
2\int_\V
\frac{\partial\etab^\top(\theta,u) }{\partial \theta_j}
\Sigmab^{-1}(\theta)
\frac{\partial\etab(\theta,u)}{\partial \theta_i}\diff \mu(u)
\\&
-2\int_\V
(\yb(u)-\etab(\theta,u))^\top 
\Sigmab^{-1}(\theta)
\begin{aligned}[t]
&\displaystyle
\left[\frac{\partial \Sigmab}{\partial\theta_j}\Sigmab^{-1}(\theta)
\frac{\partial\etab(\theta,u)}{\partial \theta_i}\right.
\\
&\displaystyle
\left.\quad
+
\frac{\partial^2\etab(\theta,u)}{\partial \theta_i\partial \theta_j}
+
\frac{\partial \Sigmab}{\partial\theta_i}\Sigmab^{-1}(\theta)
\frac{\partial\etab(\theta,u)}{\partial \theta_j}
\right]\diff \mu(u)
\end{aligned}
\\&
+
\Tr\left(\Sigmab^{-1}(\theta) \frac{\partial \Sigmab}{\partial\theta_i}\Sigmab^{-1}(\theta) \frac{\partial \Sigmab}{\partial\theta_j}\right)
+
\Tr\left(\Sigmab^{-1}(\theta) \frac{\partial^2 \Sigmab}{\partial\theta_i\partial\theta_j}\right)
\\
=&\;
2 \Mb_{ij}(\theta,\mu)
+
\Tr\left(\Sigmab^{-1}(\theta) \frac{\partial^2 \Sigmab}{\partial\theta_i\partial\theta_j}\right)
\\
&
-2\int_\V
(\yb(u)-\etab(\theta,u))^\top 
\Sigmab^{-1}(\theta)
\begin{aligned}[t]
&\displaystyle
\left[\frac{\partial \Sigmab}{\partial\theta_j}\Sigmab^{-1}(\theta)
\frac{\partial\etab(\theta,u)}{\partial \theta_i}\right.
\\
&\displaystyle
\left.\quad
+
\frac{\partial^2\etab(\theta,u)}{\partial \theta_i\partial \theta_j}
+
\frac{\partial \Sigmab}{\partial\theta_i}\Sigmab^{-1}(\theta)
\frac{\partial\etab(\theta,u)}{\partial \theta_j}
\right]\diff \mu(u).
\end{aligned}
\end{aligned}
\end{equation}





\end{document}